\documentclass[11pt]{amsart}
\usepackage[margin=1in]{geometry}
\usepackage{amsmath,amssymb,array,booktabs,longtable}
\usepackage{xurl}
\usepackage[hidelinks,hypertexnames=false]{hyperref}
\newcommand{\Z}{\mathbb Z}
\newcommand{\N}{\mathbb N}
\newcommand{\GC}{\operatorname{GC}}
\newcommand{\record}[1]{\nolinkurl{#1}}
\newcommand{\claimid}[1]{\par\nopagebreak\noindent\textit{Machine-readable claim:} \record{#1}.\par}
\newtheorem{alignedtheorem}{Theorem}
\renewcommand{\thealignedtheorem}{1a}
\newtheorem{catalogcor}{Corollary}
\renewcommand{\thecatalogcor}{1b}
\newtheorem{proposition}{Proposition}
\newtheorem{lemma}{Lemma}
\theoremstyle{definition}\newtheorem{definition}{Definition}
\setlength{\emergencystretch}{3em}
\raggedbottom
\title[Affine fullerene families and eventual periodicity]{Eventual periodicity of linear generalized Stone--Wales paths in affine fullerene families with convex lattice-face decompositions}
\author{Tatsuhiko Takimoto}
\address{Independent researcher}
\date{Version 6.3, 22 October 2026}
\begin{document}
\begin{abstract}
We study literal linear generalized Stone--Wales paths in fullerene
duals. For an integer affine family with a fixed decomposition into
embedded convex lattice polygons in the six triangular-lattice
directions, we give a uniform reduction of literal-path existence to a
fixed monadic second-order sentence over unary finite words. The
construction uses maximal face passages, their two-row outputs and all
vertex collisions. Eventual periodicity follows from the known equivalence
between MSO and regular languages. The result is a geometric definability
reduction. We apply it, together with separate arithmetic arguments,
to eleven specified catalogue families. The specification-level
periodicity theorem is Lean-proved, with a prose correspondence lemma
connecting its assumptions to the geometric hypotheses. Family-specific
Lean proofs establish the K3, C1 and D12 path classifications, O2
absence on its stated domain, and the E10 classification for $r\ge2$.
Positive-ray actual-path periodicity is Lean-proved for K3, C1, O2,
D12, D14 and E10, including the E10 boundary $r=1$. We also describe closed conditions, exact decisions on
two example rays for every ray parameter, finite checks, and a
fixed-oblique-direction extension as a separate proof sketch.
\end{abstract}
\maketitle

\section{Introduction}
A literal linear generalized Stone--Wales path is a simple alternating
triangle strip between degree-five vertices of a fullerene dual
\cite{BBS}. In an affine family, both the graph size and the number of
face passages can grow with the parameter. The question is whether
path existence nevertheless becomes periodic on an integer ray.

Our main construction represents passages and every possible
surface-vertex collision in one fixed finite-word MSO sentence.
Theorem~\ref{thm:aligned} treats embedded convex lattice faces with
fixed affine gluing. Corollary~\ref{cor:catalog} applies to eleven
specified families from Graver's catalogue \cite{Graver}; A3 uses its
separate classification and arithmetic argument \cite{A3}.
Section~\ref{sec:novelty} places this geometric reduction beside the
existing logical and fixed-boundary periodicity theorems.

\section{What is proved, and how}\label{sec:guarantees}
\begin{center}\small
\begin{tabular}{p{.23\textwidth}p{.27\textwidth}p{.38\textwidth}}
\toprule Claim & Present proof basis & Scope\\\midrule
Theorem~\ref{thm:aligned} & Lean specification theorem; prose geometric correspondence & Convex six-direction disk decompositions, with fixed affine gluing and eventual validity; fixed-residue variants.\\
Logical periodicity step & External MSO/regular-language theorem & One fixed sentence over unary finite words.\\
Corollary~\ref{cor:catalog} & Geometric decompositions and arithmetic; direct Lean family proofs & Eleven specified valid positive-parameter families.\\
K3, C1, O2, D12, D14, E10 & Lean quotient-graph proofs & K3 iff for $r,s\ge2$; C1 iff for $a,b\ge1$; D12 iff for $r,s\ge1$; O2 absence for $t,d\ge2$; E10 iff for $p,s\ge1,r\ge2$; all six positive-natural ray results, including E10 $r=1$.\\
Family boundaries & Lean explicit paths and finite checks & O2 paths on $(1,d)$, $d\ge3$, and $(t,1)$, $t\ge3$; D14 paths on $(1,s)$, $s\ge2$, and $(r,1)$, $r\ge2$. Small remaining boundary decisions are finite evidence.\\
Other closed conditions and examples & Prose geometry, Lean predicate arithmetic, explicit witnesses and checkers & The domains and finite ranges specified in Sections~\ref{sec:conditions}, \ref{sec:arithmetic} and \ref{sec:finite}; C1/O2 IPR is prose.\\
% D12-SCOPE-BEGIN
D12 classification, further proofs & Two prose/computer-assisted proofs, in addition to the Lean proof & Positive integer $r,s$, with the quotient, geometric and certificate semantics stated in Section~\ref{sec:d12-classification}.\\
% D12-SCOPE-END
Fixed-oblique extension & Proof sketch & The four transfer obligations in Section~\ref{sec:extension}.\\
\bottomrule
\end{tabular}
\end{center}
The formal family claims concern the explicit quotient models; their
catalogue identification uses the stated coordinate conventions.
The convex-face result has two separately developed derivations on
the common six-direction domain. Human proof peer review and the
priority comparison remain open. The supplementary
\record{VERIFICATION-RECORD.md} records provenance, implementation
identifiers, independence qualifications and validation history.

\section{Definitions, signatures and valid rays}\label{sec:definitions}
\subsection{The literal path}
\begin{definition}\label{def:literal}
In the dual $G$ of a fullerene, a literal linear gSW path is a sequence
$p_0,\ldots,p_e$ with odd $e\ge3$, all vertices distinct, consecutive
and distance-two vertices adjacent, $\deg p_0=\deg p_e=5$, and
$\deg p_1=\deg p_{e-1}=6$.
\end{definition}
\claimid{def:literal}
This is Definition 1(a) of \cite{BBS}. Internal degree-five vertices
are allowed. Even vertex count is part of the definition, not a new
obstruction. Positive means existence of such a path; no chemical
barrier, radial operation, or operation permitting heptagons is inferred.
A fullerene is a simple three-connected cubic plane graph with only
pentagonal and hexagonal faces. Its dual is a spherical triangulation
with twelve degree-five vertices and all other degrees six.

When every three-cycle is facial and an edge has just the two common
neighbours supplied by its incident triangles, an oriented triangle
$(a,b,c)$ has a unique next vertex: the common neighbour of $b,c$
different from $a$. A simple path is a prefix of this forced sequence.
If $\rho$ is its first repeated-vertex index, a legal endpoint must
occur strictly before $\rho$, not at $\rho$.
\label{prop:forced}\claimid{prop:forced}
For fullerene duals, cyclic edge-connectivity
\cite{Connectivity} excludes nonfacial three-cycles and supplies the
two-common-neighbour property. The finite-graph decision argument
applies once this local geometric hypothesis and the start are fixed.

\subsection{Signatures and eventual validity}\label{def:signature}
\claimid{def:signature}
Work in the triangular lattice $\Lambda=\Z^2$, with
$R(x,y)=(-y,x+y)$. A signature specifies a fixed finite set of polygon
charts, paired boundary segments, corner equivalence classes, and
transition maps $z\mapsto R^jz+d$. Corner coordinates and translations
are integer affine functions of the parameters; rotations and gluing
combinatorics are fixed. Along a ray $p(t)=at+b$, $t\in\N$, we require
eventually positive edge lengths, embedded chart interiors, only the
specified seam identifications, noncrossing macro edges, and a valid
simple fullerene dual. Every cone has a marked corner label. Seams
need an ownership convention, and full corner equivalence classes are
used. Positive parameters and positive area alone do not prove validity.

A component $a_it+b_i$ eventually lies in the positive integers exactly
when $a_i>0$ or $a_i=0,b_i\ge1$. In the first case an entry bound is
$\max(0,\lceil(1-b_i)/a_i\rceil)$. Replacing 1 by 2 gives the analogous
bound for a lower limit two. A negative slope has no infinite positive
tail. These arithmetic entry bounds do not certify the quotient.

\section{The main result: convex lattice-face decompositions}\label{sec:main}
\begin{alignedtheorem}[Convex lattice-face rays]\label{thm:aligned}
For all sufficiently large integers $t$, suppose a valid simple
fullerene dual has a decomposition into a fixed finite number of
embedded closed convex lattice-polygon disks. Every side has positive
length and one of the six lattice directions. Straight angles and
auxiliary degree-six corners are permitted. No distinct boundary
points of one disk are identified. Corner coordinates and transition
translations are integer affine in $t$, while rotations, side pairings
and corner identifications are fixed. All degree-five vertices occur
among the marked corners, with a fixed label set. Then the set of
$t$ for which a literal path exists is eventually periodic.
The conclusion also holds when these hypotheses hold separately on
finitely many residue classes modulo a fixed integer.
\end{alignedtheorem}
\claimid{thm:aligned}
Thus some $t_0\ge0,P\ge1$ satisfy the same decision at $t$ and $t+P$
for all $t\ge t_0$. This is an existence theorem for a threshold and
positive period. The following correspondence lemma connects the
geometric hypotheses to the specification-level theorem, while
Section~\ref{sec:proof} gives the passage-to-word proof.

\paragraph{Prose correspondence lemma.}
The geometric hypotheses of Theorem~\ref{thm:aligned} determine a fixed
\record{Sig} $S$ and a tail on which
\record{Valid} $S\,t$ and \record{TrianglesInCharts} $S\,t$ hold;
the quotient \record{graph} $S\,t$ is the stated fullerene dual and
\record{LiteralPath} is Definition~\ref{def:literal}.

\emph{Prose justification.}
For a convex lattice polygon with the six permitted side directions,
its supporting lines give exactly three closed strips
$L_x\le x\le U_x$, $L_y\le y\le U_y$,
$L_s\le x+y\le U_s$. Their extrema are selected from finitely many
affine corner coordinates; after their orders stabilize, all six
bounds are integer affine. Retain the fixed charts, rotation exponents,
affine translations and marked corners as the fields of $S$. For each
paired side, use its prescribed transition as a gluing record.

\emph{No extra identifications from unrestricted records.}
Let $g$ unfold a paired segment $A\subset K$ onto $g(A)\subset K'$,
and let $\ell$ be its supporting line. The consistently oriented disks
$g(K)$ and $K'$ occupy opposite closed supporting halfplanes, so
$g(K)\cap K'\subset\ell$. Their overlap there is an interval $I$
containing $g(A)$. If $I$ continues beyond an endpoint of $g(A)$,
both supporting edges continue straight at that point. The two disks
each contribute angle $\pi$. A fullerene dual has total cone angle
$5\pi/3$ or $2\pi$, so the point is flat, the two disks exhaust its
neighbourhood, and their actual seam must continue across it. The
orientation-preserving lattice isometry along the continued seam is
the same $g$: it agrees on a nontrivial line segment. Repeat across
the finitely many side subdivisions until reaching the endpoint of
$I$. Thus every pair identified by this unrestricted record is truly
the same surface point, even if it lies beyond the named subdivision.
Conversely, the prescribed paired sides are included in the records.
A self-pairing of distinct points is excluded by the embedded-closed-disk
hypothesis; an identity record is harmless. Transitive closure therefore
gives exactly the prescribed seam and full corner equivalences, without
extra identifications. This supporting-line and cone-angle argument is
a prose geometric argument. Thus sites
modulo these records are precisely the surface vertices, and common-chart
unit steps are precisely the graph edges. Closed-disk embeddedness
supplies the injectivity clause of \record{Valid}; the fullerene degrees
and marked-corner hypothesis supply its degree clauses.

Use the same external fullerene fact as in Section~\ref{sec:definitions}:
all dual three-cycles are facial, by cyclic edge-connectivity
\cite{Connectivity}. Each edge therefore has exactly the two distinct
opposite vertices of its incident triangular faces as common neighbours.
This gives \record{TwoCommonNeighbors}. Since chart boundaries follow
lattice edges, no open unit triangle is cut by a boundary; the covering
decomposition puts every triangular face inside a closed chart.
Every graph triangle is facial, so its three vertices lift there with
all three differences unit vectors. This is exactly
\record{TrianglesInCharts}. These geometric identifications and the external fullerene fact
supply the geometric correspondence.

Apply the specification-level periodicity theorem on the resulting tail. For fixed
residue branches write $t=r+qn$ and apply the construction to each
branch. The residue-class theorem combines the branches in the
original parameter, assuming eventual
agreement of the decision with each branch's \record{PathExists}.
That agreement is supplied by this prose representation lemma;
the combination of eventual periods is Lean-proved.
\begin{catalogcor}[Specified catalogue families]\label{cor:catalog}
On every integer affine ray eventually remaining in the valid positive
integer domain of each of A1, A2, B2, C1, K3, E10, D12, D14, D1 and O2,
literal-path existence is eventually periodic. The same conclusion
holds for A3 on rays eventually satisfying $k>l>0$, by its separate
classification and arithmetic argument~\cite{A3}.
\end{catalogcor}
\claimid{cor:catalog}
There are eleven families in this statement. Positive boundaries with
a parameter fixed at 1 are included, notably K3 and O2 and E10 with
$r=1$. Collapsed zero boundaries are not. The corollary gives no
two-variable closed classification merely from ray periodicity.
The K3, C1, O2, D12 and D14 parts have direct family-specific ray proofs
on their quotient graphs, including positive boundary parameters. E10 also
has a formal actual-path ray proof: its $r=1$ branch applies the formal
general periodicity theorem with the validity and triangle-lifting
hypotheses proved on the E10 model. These six results use no general
prose correspondence lemma. Their ray declarations start with
positive natural coordinates and nonnegative slopes; an eventually positive
integer ray is covered by reindexing its positive tail. The other family
routes retain their stated scope.
A3 generally has changing parent-edge directions and is not silently
placed under Theorem~\ref{thm:aligned}. The catalogue terminology and
corrections are Graver's \cite{Graver}; the constructions are not claimed
as new. Additional leapfrog families are discussed only as prospects
in Section~\ref{sec:extension}.

\section{Proof of the convex-face theorem}\label{sec:proof}
The proof uses the common six-direction domain and its restricted
arithmetic-to-word reduction. The fixed-oblique-direction sketch is
treated separately in Section~\ref{sec:extension}.

\subsection{Maximal face passages and their outputs}\label{lem:rows}
\claimid{lem:rows}
Rotate a passage direction to $(1,0)$. A suitable phase of the forced
vertex sequence is
\begin{equation}\label{eq:rows}
 q_{2j}=(j,h),\qquad q_{2j+1}=(j,h+1).
\end{equation}
A convex six-direction polygon is specified by
\[
 L_x\le x\le U_x,\qquad L_y\le y\le U_y,\qquad
 L_s\le x+y\le U_s.
\]
Its bounds are affine in $t$ after the finitely many max/min orders
stabilize. Redundant collinear sides do not alter this description.
With $L_y\le h\le U_y-1$, the triangle
$T_i=\{q_i,q_{i+1},q_{i+2}\}$ lies in the polygon precisely in the ranges
\begin{equation}\label{eq:triangle-bounds}
\begin{aligned}
 A_e&=\max(L_x,L_s-h),& B_e&=\min(U_x-1,U_s-h-1),\\
 A_o&=\max(L_x,L_s-h-1),& B_o&=\min(U_x-1,U_s-h-2).
\end{aligned}
\end{equation}
Even indices are $2j$ for $A_e\le j\le B_e$, odd indices are $2j+1$
for $A_o\le j\le B_o$. Put $\delta=A_e-A_o$ and
$\epsilon=B_o-B_e+1$, both in $\{0,1\}$. The full triangle-index set
is empty or the interval
\begin{equation}\label{eq:passage}
 \ell=2A_e-\delta\le i\le u=2B_e+\epsilon.
\end{equation}
This includes the one-odd-triangle case $A_e=B_e+1$,
$\delta=\epsilon=1$. It must not be discarded as an empty even range.

The first two vertices of an entered passage already belong to the
previous passage. Its new output is only
$q_{\ell+2},\ldots,q_{u+2}$, comprising
\begin{equation}\label{eq:output}
\begin{array}{ll}
 y=h:   &A_e+1\le x\le B_e+1,\\
 y=h+1: &A_e+1-\delta\le x\le B_e+\epsilon.
\end{array}
\end{equation}
For example the upper row has
$x\ge L_x+1-\delta$, $x+y\ge L_s+2-\delta$,
$x\le U_x-1+\epsilon$, and $x+y\le U_s+\epsilon$.
Empty rows are allowed, but a nonempty passage has nonempty output.

Root and terminal passages can be truncated. There are finitely many
marked cone corners and at most ten oriented initial triangles at each.
The initial triangle fixes $h$ and an index of the form twice an affine
expression plus fixed parity. Store its first two vertices separately.
An endpoint corner and choice of output row similarly determine the
last index. Truncation adds affine bounds on $x$. Both cuts are made
when the whole path lies in one passage. Internal degree-five corners
are not compulsory stopping points. The two endpoint neighbours must
be distinct from all degree-five labels.

\subsection{States, transitions and all collisions}\label{lem:restricted}
\claimid{lem:restricted}
A state is $(\tau,h)$, where the finite type $\tau$ records face,
direction, max/min branches, parity, subdivided entry/exit side, and
root/terminal information. A constant $C$ bounds $|h|\le Ct+C$.
The number of states or face passages need not be bounded independently
of $t$. The relevant relations are Boolean combinations of
\begin{equation}\label{eq:restricted}
 \pm h\le at+b,\qquad \pm h\pm h'\le at+b,
 \qquad at+b\ge0,
\end{equation}
and fixed-modulus congruences. A relation such as $2h+h'\le at+b$
with distinct free variables is not part of this language. Substituting
one variable for another can be reduced by fixed parity and rounding.

Within a fixed type, each endpoint's three normal coordinates has the
form $c(t)$, $h+c(t)$ or $-h+c(t)$. When the last ordered pair of a
passage is glued to the next one's first ordered pair, the directions
are rotated together. In direction-aligned coordinates the linear part
of this transition is the identity, so $h'=h+c(t)$, not a doubling map.
The paired unit edge determines the opposite face; its interior
contains no macro corner. The conditions therefore have
form~\eqref{eq:restricted}.

Each output row fixes one of $x,y,x+y$ at $\pm h+c(t)$ and imposes
affine bounds in all three normals. For distinct row normals the
determinant is $\pm1$: their intersection is integral, and the third
normal is their sum or difference. Substitution into both rows gives
\eqref{eq:restricted}. Equal normals first require equal values;
the remaining integer interval is nonempty exactly when every lower
bound is at most every upper bound. The same restricted form results.
Between different faces, compare through every shared side and every
common marked corner label. Rotations permute the three normals with
signs. Embeddedness and fixed gluing ensure these exhaust quotient
identifications, even when faces share several sides. Auxiliary flat
corners are included. Comparison with the root's two saved vertices
uses the same formulas. Since an embedded face row is injective,
pairwise disjoint outputs and their disjointness from the root are
exactly the full vertex-distinctness condition.

\subsection{Folding the restricted arithmetic into a finite word}\label{lem:folding}
\claimid{lem:folding}
For $t\ge2$, write a height variable uniquely as $h=jt+y$, $0\le y<t$.
The quotient $j$ ranges over a fixed finite set. Set $t=2n+\eta$,
$\eta\in\{0,1\}$, and split $y$ into
\begin{equation}\label{eq:fold}
 L:y=v,\qquad R:y=t-1-v\quad(0\le v<n),\qquad
 M:y=n\quad(\eta=1).
\end{equation}
Original type, $j$ and $L/R/M$ are finitely many copies of the positions
of a unary word of length $n$. A middle copy is placed at the first
position and interpreted as coordinate $n$, independent of that position.

For two ordinary positions the inequalities become
\[
 \alpha v+\beta w\le 2cn+d,\qquad \alpha,\beta\in\{-1,1\}.
\]
The coefficient of $n$ is even. If $\alpha=-\beta$, then for $c\ne0$
truth stabilizes, while $c=0$ is order with a fixed successor offset.
If $\alpha=\beta=1$, the sum lies between 0 and $2n-2$. The only
nontrivial cases are $c=0$, a finite list near the left end, and $c=1$.
For $c=1$, putting $a=n-1-v,b=n-1-w$ gives
$a+b\ge-d-2$, either automatic or the exclusion of finitely many
near-end pairs. The case $\alpha=\beta=-1$ is its integer complement.
If one position is a middle copy, only one variable remains and the
nontrivial inequalities are fixed distances from an end. Two middle
copies give an affine truth value that stabilizes.

Fixed-modulus congruences are represented by cyclic MSO colours on
the word: the first position has colour zero and each successor
increments the colour. The final colour determines $n$ modulo the
fixed modulus. Finite copies of element or set quantifiers are expanded
copy by copy. Thus every fixed Boolean combination used above is
MSO-definable. This argument does not permit an arbitrary Presburger
relation such as $v+w\le n$; it uses the even coefficient obtained by
the folding step.

\subsection{The exact path sentence and unary regularity}\label{lem:mso-path}
\claimid{lem:mso-path}
Choose a nonempty state set $S$, root $a$ and terminal $b$. Terminal
states have no outgoing transition and root states no incoming one.
Record the parity of the number of vertices already output as a finite
type. Initially it is even, since the root has two saved vertices;
each transition adds the parity of $u-\ell+1$.
Impose the following MSO conditions:
\begin{enumerate}
\item If $a=b$, $S=\{a\}$. Otherwise $a$ has in/out degrees $(0,1)$,
      $b$ has $(1,0)$, and every other state has $(1,1)$ within $S$.
\item Every nonempty proper subset of $S$ has a transition, in either
      direction, across its boundary. This is unary-set connectivity.
\item Outputs of distinct states are disjoint, and all outputs are
      disjoint from the two saved root vertices.
\item Endpoint degrees are $5,6,6,5$ as in Definition~\ref{def:literal}.
      Total vertex count is even and at least four.
\end{enumerate}
For one state, the last size condition is $u-\ell+1\ge2$; with two
or more nonempty-output states it is automatic. Degree conditions are
membership or nonmembership in the finite marked-cone set. No binary
relation quantifier is used.

On a finite set, the first two conditions give one directed path
through all of $S$: a reached cycle would create an extra incoming edge,
and a separate cycle would contradict connectivity. Join the passages
in this order. Ordered seam pairs give consecutive and skip adjacency;
the third condition gives distinct vertices. Size, parity and degree
conditions give exactly the literal path.

Conversely, split any literal path into maximal consecutive passages
in its convex faces, truncating only at its ends. Equations
\eqref{eq:triangle-bounds}--\eqref{eq:output} recover their states.
Nonempty output and path simplicity prevent a repeated state. Any extra
eligible successor within $S$ would share the next output vertex with
the actual successor, contradicting disjointness. Hence the selected
states satisfy all four conditions. One-passage paths and paths with
internal cones are included in both directions.

Finite-word MSO defines exactly the regular languages
\cite[Section 3.1, Theorem 3.2]{MSO}. Apply this established result
only at this last logical step. For each $\eta$, a deterministic
automaton over one letter eventually repeats a state, so its accepted
lengths $n$ are eventually periodic. Combine $t=2n+\eta$ and, if
necessary, the initially fixed residue branches. This proves
Theorem~\ref{thm:aligned}. The geometric passage construction supplies
the fixed sentence to which the external automata theorem applies.

\section{Applying the theorem to the specified families}\label{sec:families}
\label{prop:family-decompositions}\claimid{prop:family-decompositions}
D12 uses two subdivided hexagons and D14
two convex hexagons whose long sides are subdivided with lengths
$r,s,r$. O2 uses two hexagons and six triangles. K3 uses two caps,
each divided into four triangles, and three hexagons. These
decompositions retain the positive parameter-1 boundaries.

D1's original thirty-degree sides do not directly satisfy the six
direction hypothesis. Cut a triangle $A_iB_iW_i$ from each oblique side
and move it to the back by the same lattice isometry. For example,
\[
 A_0=(r,0),\quad B_0=(r+p,p),\quad W_0=(r,p),\quad C=(-p,r+2p),
\]
and $W_0-C=\alpha(A_0-C)+\alpha(B_0-C)$ with
$\alpha=(r+p)/(2r+3p)$ and $0<2\alpha<1$. The pieces are contained
and disjoint. The result is eight convex pieces: one regular hexagon,
six quadrilaterals, and the back hexagon. New corners are flat degree
six. Area equality alone is not used to infer the same quotient.

For E10 put $m=2r+s$, tube circumference $C=3m$. Move the longitudinal
cut from $(p,p)$ to $(0,2p)$, obtaining three lattice parallelograms
of width $m<C$ and two caps: five disks, also when $r=1$. Cut-corner
positions involve $p\bmod C$. If $C$ grows, $\lfloor p/C\rfloor$
eventually stabilizes; if $C$ is constant, use fixed residue classes.
Delete duplicate auxiliary corners and fix their order to obtain
affine coordinates and fixed gluing. This does not classify E10 with
$r=1$.

\begin{center}\small
\begin{tabular}{p{.19\textwidth}p{.69\textwidth}}\toprule
Family and domain & Reason for all valid affine rays\\\midrule
A1, $r\ge1$ & Integer GC absence, including the separate $r=1$ case; eventual period one.\\
A2, $p\ge1$ & Nonconstant rays eventually give nonprimitive GC images; constant rays are constant graphs.\\
B2, $p,r\ge1$ & Its closed condition is a finite disjunction of affine equalities.\\
C1, $a,b\ge1$ & Lean graph iff $b=1$ and actual-path ray periodicity; no general prose bridge.\\
A3, $k>l>0$ & Separate complete classification and arithmetic; Section~\ref{sec:arithmetic}.\\
K3, $r,s\ge1$ & Lean literal-path ray periodicity, including parameter 1; no general prose bridge is used.\\
E10, $p,r,s\ge1$ & Lean actual-path ray periodicity; for $r=1$, the formal general theorem applies with model hypotheses proved in Lean.\\
D12, $r,s\ge1$ & Lean actual-path ray periodicity, including parameter 1; no general prose bridge.\\
D1 & The convex decomposition on valid positive integer parameters.\\
D14, $r,s\ge1$ & Lean actual-path ray periodicity, including parameter 1; no general prose bridge.\\
O2, $t,d\ge1$ & Lean actual-path ray periodicity, including parameter 1; Lean absence for $t,d\ge2$. No general prose bridge.\\\bottomrule
\end{tabular}
\end{center}
These decompositions and the separate arithmetic results prove
Corollary~\ref{cor:catalog} on its stated domains.

\section{Closed conditions and the scope of partial results}\label{sec:conditions}
The following closed conditions concern valid catalogue graphs on
the displayed domains.

\subsection{The specified closed domains}\label{prop:closed}
\claimid{prop:closed}
\begin{center}\small
\begin{tabular}{p{.19\textwidth}p{.36\textwidth}p{.31\textwidth}}\toprule
Family/domain & Literal path condition & Proof and domain\\\midrule
A1, $r\ge1$ & Never positive & Integer GC for $r\ge2$; at 1 no degree-six vertex.\\
A2, $p\ge1$ & $p=1$ & $C_{60}$ witness and integer GC.\\
C1, $a,b\ge1$ & $b=1$; IPR iff $a,b\ge2$ & Graph iff in Lean; IPR equivalence by prose geometry.\\
B2, $p,r\ge1$ & $p=1$ or $|p-r|=1$ & Prose proof on the whole positive domain.\\
K3, $r,s\ge2$ & $s\equiv\pm1\pmod r$ & Graph iff in Lean on the displayed domain.\\
O2, $t,d\ge2$ & Never positive & Graph absence in Lean; IPR by prose geometry.\\
E10, $p,s\ge1,r\ge2$ & $m=2r+s$ and $p\equiv d\pm1\pmod m$ for some $d\in\{0,r,r+s\}$ & Graph iff in Lean on the displayed domain; the ray theorem also covers $r=1$.\\\bottomrule
\end{tabular}
\end{center}
K3 writes $s=qr+u$: $u=0$ is excluded by integer GC,
$2\le u\le r-2$ by belts of widths $r,u,r-u$, and $u=1,r-1$
by explicit simple paths. The overlapping small cases $r=2,3$ are
separate. The iff is not extended to $r=1$ or $s=1$. E10's candidate
congruences come from the positions of corners around the tube;
reaching a cone does not alone establish simplicity. The E10 argument uses eighteen parameter regions, their walls and
explicit witnesses.

\subsection{Partial branches and boundary examples}\label{prop:partial}
\claimid{prop:partial}
For O2, absence holds on $t,d\ge2$, with IPR on that domain supplied
by prose geometry. Boundary paths exist at O2$(1,d)$ for $d\ge3$ and
O2$(t,1)$ for $t\ge3$. At O2$(1,1)$, O2$(1,2)$ and O2$(2,1)$,
the path decisions have finite computational evidence. Every constant
graph has a constant ray decision, independently of computing it.
For D14, absence holds for $r,s\ge2$, and explicit paths exist at
$(1,s)$ for $s\ge2$ and $(r,1)$ for $r\ge2$. These are Lean graph
lemmas; the D14 ray theorem includes all positive boundary rays.
D14$(1,1)$ has a finite path witness. For D1 with both parameters at
least two, a belt-based absence argument is a proof sketch. Common
integer scaling by at least two gives absence in the specified
leapfrog domains, without classifying their primitive points.

For E10$(p,1,s)$, $m=s+2$, explicit witnesses give
sufficiency when $p<m$ or
$p\bmod m\in\{0,1,2,m-2,m-1\}$. General necessity is not proved.
Among the 84 positive $(p,s)$ with
$N=12+2s^2+12s+12p(2+s)\le740$, exactly E10$(9,1,4)=C_{740}$
is negative; all 83 smaller examples are positive. E10$(17,1,11)$
and $(31,1,7)$ give further negatives at $C_{3038},C_{3542}$.
These three examples are non-IPR. Exhaustive construction and
independent checks establish the stated minimality inside the $r=1$
E10 family.

For D12, the following partial branches apply:
$\gcd(r,s)\ge2$ gives absence; $r\ge3$, $s=qr\pm1$, $q\ge1$
gives existence by the separately derived simple witnesses. The
additional $|r-3s|=1$ and $r\ge4s$ branches are construction and
absence arguments with their stated geometric conditions. The full D12 classification, with a Lean proof, is given in Section~\ref{sec:d12-classification}.

\section{Arithmetic on affine rays and exact examples}\label{sec:arithmetic}
\subsection{The determinant reduction}\label{lem:divisibility}
\claimid{lem:divisibility}
Let $A(t)=\alpha t+\beta$ be eventually positive and
$B(t)=\gamma t+\delta$. A constant $A$ gives a fixed-modulus
congruence. If $\alpha>0$, then
\begin{equation}\label{eq:determinant}
 A\mid B\quad\Longrightarrow\quad A\mid\alpha\delta-\gamma\beta.
\end{equation}
For nonzero right side, divisibility is impossible once
$A>|\alpha\delta-\gamma\beta|$. For zero right side, the ratio
$B/A=\gamma/\alpha$ is constant and is either always integral or
never integral. Finite Boolean combinations of affine divisibilities
are eventually periodic; a growing modulus gives eventual constancy.
Take a common multiple of fixed moduli and a maximum stabilization
threshold. This elementary arithmetic is not proposed as new number
theory. The following exact predicates have Lean arithmetic proofs.

For all integer ray coefficients, define parameters by
$(at+b).\operatorname{toNat}$. The formal statements for all rays use
\emph{domain and positive predicate}; they do not call out-of-domain
graphs positive. On an eventually positive domain the conversion
agrees with the integer affine parameter.
\label{thm:k3-arithmetic}\claimid{thm:k3-arithmetic}
For K3, eventual periodicity holds for
$\mathrm{K3Dom}\wedge\mathrm{K3Pos}$ on all integer affine rays,
where the domain is $r,s\ge2$ and the predicate is
$r\mid(s-1)$ or $r\mid(s+1)$. If $r$ is nonconstant on a valid tail,
the predicate is eventually constant. If $r=b\ge2$ is
constant, period $b$ is valid from any specified
domain-entry threshold. These need not be least periods.
\label{thm:e10-arithmetic}\claimid{thm:e10-arithmetic}
For E10, the corresponding all-ray statement holds
for domain $p,s\ge1,r\ge2$ and the six divisibilities in the E10
condition. If $2r+s$ is nonconstant, the predicate is eventually
constant; when it is constant, that
positive modulus is a period on the given domain tail.
\label{thm:other-arithmetic}\claimid{thm:other-arithmetic}
For B2, C1, A1 and A2, eventual constancy holds for the domain-qualified
predicates. The conjunction/disjunction/negation closure and domain
qualification are proved in the same arithmetic layer.

\subsection{A3 is a separate variable-direction case}\label{thm:a3-arithmetic}
\claimid{thm:a3-arithmetic}
For $k(t)=at+b,l(t)=ct+d$ set $\Delta=ad-bc$. If $\Delta\ne0$,
$\gcd(k,l)=\gcd(k,l,\Delta)$, so primitivity is periodic modulo
$|\Delta|$. When $\Delta=0$ and the ray is nonconstant, it is
eventually an increasing integer multiple of a fixed primitive
direction and hence negative by integer GC. Constant rays are constant.
In the primitive case put $m=\min(l,k-l)$. The A3 criterion is
$m=1$, or quotient $q=\lfloor k/m\rfloor\ge3$ with $m\mid(k\pm1)$,
or $q=2$ with $r=k-2m=1$ or $r\mid(m\pm1)$. Branches and quotient
stabilize, including an integer limiting ratio by checking the sign
of the constant remainder. The persistent $q=2,r=0$ growing branch
is not primitive. Apply~\eqref{eq:determinant} to the remaining tests.

The first paper proves the full graph classification and a graph-path
ray theorem. Its arithmetic proof establishes a threshold and positive
period; an independent arithmetic derivation permits period $|\Delta|$
for nonconstant $\Delta\ne0$ rays and period 1 otherwise. These are
valid periods, with no claim of leastness. After classification,
eventual periodicity is also a consequence of
\cite[Theorem 1.15]{BGW}.

\subsection{Four previously solved rays}\label{ex:exact-rays}
\claimid{ex:exact-rays}
\begin{center}
\begin{tabular}{lll}\toprule
Ray, $t\ge0$ & Positive values & Valid period and threshold\\\midrule
D12$(124+3t,41+t)$ & All $t$ & $P=1,t_0=0$\\
D12$(125+3t,41+t)$ & None & $P=1,t_0=0$\\
K3$(43+2t,170+5t)$ & $40,42$ & $P=1,t_0=43$\\
E10$(200+5t,17+2t,23+t)$ & $3,10,17,34$ & $P=1,t_0=35$\\\bottomrule
\end{tabular}
\end{center}
Compressed-word calculations and independent family-condition
arguments give these exact all-parameter decisions.
\label{ex:k3-formal}\claimid{ex:k3-formal}
For every $t\in\N$, the K3 arithmetic theorem gives
$\mathrm{K3Pos}(43+2t,170+5t)\leftrightarrow(t=40\lor t=42)$.
\label{ex:e10-formal}\claimid{ex:e10-formal}
The E10 arithmetic theorem gives the corresponding predicate
iff $t\in\{3,10,17,34\}$. The determinant constants used are
$123,127$ for K3 and $142,144,739,749,681,691$ for E10. Divisibility
bounds $t$ before the remaining cases are checked by kernel arithmetic.
These are all-$t$ arithmetic proofs, not graph iff proofs by themselves.

For every $s\ge41$, D12$(3s+2,s)$ is negative and
D12$(3s+1,s)$ has a path with $56s+12$ distinct vertices. The
computer-assisted proof checks 2,058 nonnegativity conditions and
756 no-integer-solution conditions, 2,814 in total. A separate
construction at $s=58$ gives a 3,260-vertex positive witness.

\section{Specification and family-specific graph proofs}\label{sec:formal}
For every fixed signature $S$ and $t_1$, the specification-level
theorem states that if \record{Valid} $S\,t$ and
\record{TrianglesInCharts} $S\,t$ hold for every $t\ge t_1$, then
some $t_0$ and positive $P$ satisfy
\record{PathExists} $S\,t\leftrightarrow$
\record{PathExists} $S\,(t+P)$ for all $t\ge t_0$.
These assumptions require the local validity and triangle-lifting
properties; the geometric correspondence lemma supplies them from
the convex disk decomposition. The residue-class variant combines
fixed branches after their decisions agree with the original family.

\subsection{K3 graph classification and ray periodicity}\label{thm:k3-sufficient}
\claimid{thm:k3-sufficient}
For natural $r,s\ge2$, K3 has a literal path if and only if
$s\equiv\pm1\pmod r$. A classification of graph triangles in the
cap or tube proves that no edge has three common neighbours and
supplies the forced-continuation argument for necessity.

Literal-path existence is eventually periodic on every ray
$(r,s)=(r_0+at,s_0+bt)$ with $r_0,s_0\ge1$ and $a,b\ge0$.
The proof constructs the common-neighbour property directly.
Nonconstant parameter-one boundary rays use explicit eventual paths;
constant rays have constant decisions. The modular iff has domain
$r,s\ge2$.

The specification uses eleven convex polygons, equivalence closure of
explicit affine side gluings, and adjacency from unit differences inside
a common chart. The literal path and actual neighbour-set degree
 definitions are those of the first paper's specification.

The proof uses a tube construction. A glued-site invariant
is $(\mathrm{hgt},\mathrm{ang}\bmod6r)$. It is preserved by all
42 directed edge gluings. It is not globally injective: the two
central cap triangles are not separated by it. On tube points with
horizontal coordinate difference at most one, however, equality of
quotient vertices forces equality of points. The explicit paths are
\[
 P_n^+=(-(n\bmod2),\lfloor(n+1)/2\rfloor),\qquad
 P_n^-= (n\bmod2,\lfloor n/2\rfloor),
\]
with $4s-2$ vertices for $s=qr+1$ and $4s+4$ for $s=qr-1$.
Chart membership, consecutive and skip adjacency, distinctness, and
the four endpoint degree requirements are proved on the quotient.
The invariant is not used to claim a global vertex normal form.

\label{lem:k3-corners}\claimid{lem:k3-corners}
All twelve labelled corner representatives have degree five. These
representatives supply the endpoint degree-five witnesses.
\label{lem:k3-hex-interior}\claimid{lem:k3-hex-interior}
Hexagon interiors have degree six.
\label{lem:k3-tri-interior}\claimid{lem:k3-tri-interior}
The six noncentral cap-triangle interiors have degree six.
The necessity argument uses the cap/tube triangle classification
and common-neighbour bound together with these local degree facts.

\subsection{C1 graph classification and ray periodicity}\label{sec:c1-formal}
For C1, natural $r,s\ge1$ admit a literal path if and only if $s=1$.
Here $r,s$ correspond to the paper's $a,b$. Actual-path existence is
eventually periodic on every ray $(r_0+ut,s_0+vt)$ with
$r_0,s_0\ge1$ and $u,v\ge0$, including constant rays.
Eventually positive integer rays are covered by shifting their
starting index to the positive natural domain.

The eight-cell specification has four hexagons with alternating sides
$s,r,s,r,s,r$ and four triangles of side $r$. A vertex bijection preserving adjacency identifies
its quotient with the lattice surface in $\mathbb Z^4$ cut out by
$0\le x_i\le r+s$, $\sum_i x_i=2r+s$, and some $x_i\in\{0,r+s\}$.
The bijection preserves literal-path existence. The proof establishes degrees and excludes three common
neighbours of an edge directly, then uses forced strips for $s\ge2$
and an explicit path for $s=1$.

\subsection{O2 absence and ray periodicity}\label{sec:o2-formal}
For natural $t,d\ge2$, O2 has no literal path. Actual-path existence
is eventually periodic on every ray $(t_0+an,d_0+bn)$ with
$t_0,d_0\ge1$ and $a,b\ge0$, including parameter-one and constant
rays. Eventually positive integer rays are reindexed to their
positive tail.

The eight-cell specification uses two hexagons with alternating sides
$2t,d,2t,d,2t,d$, each long side split at its midpoint, and six triangles
of side $t$. A graph equivalence identifies it
with two overlaid hexagons sewn along the three $d$-sides, with a rhombic
pocket inserted into each $2t$-side slit. The proof directly rules out
three common neighbours on an edge, locates degree-five vertices at the
twelve corners, and covers the 120 starting triangles by ten closed
strips.

Explicit boundary paths are proved for $(1,d)$ with $d\ge3$ and
$(t,1)$ with $t\ge3$. Path existence at $(1,1),(1,2),(2,1)$ is supported by finite checks. These three constant graphs do not
require their existence decision for the ray theorem: any fixed graph
has a constant decision. Explicit finite witnesses give paths there.

\subsection{D14 ray periodicity and auxiliary boundary lemmas}\label{sec:d14-formal}
For D14, actual-path existence is eventually periodic on all natural
rays $(r_0+an,s_0+bn)$ with $r_0,s_0\ge1$ and nonnegative slopes,
including parameter-one and constant rays. Eventually positive
integer rays are reindexed to a positive natural tail.

The specification consists of two hexagons with alternating side lengths
$s,2r+s$, subdividing each long side as $r,s,r$. A graph
equivalence maps the quotient to two lattice hexagons
whose boundary circles are identified with an offset between their
corners. A seam vertex is recorded by its arclength modulo $6r+6s$.
The proof directly excludes three common neighbours on an edge and,
for $r,s\ge2$, locates degree-five vertices at the twelve corners.
Twelve closed strips cover all 120 starts.

Absence holds for $r,s\ge2$. Explicit paths exist at $(1,s)$ for
$s\ge2$ and $(r,1)$ for $r\ge2$. At $(1,1)$, a finite path witness
supplies the decision; the constant-ray statement follows directly
from constancy of the graph.

\subsection{E10 graph classification and all positive rays}\label{sec:e10-formal}
For natural $p,s\ge1$ and $r\ge2$, the E10 quotient admits a literal
path exactly under the six congruences in the closed-condition table.
An isomorphism identifies the five-polygon quotient with two hexagonal
caps and a tube of circumference $3m$ and height $2p$, where $m=2r+s$
and the twist is $p\bmod3m$. Degree five occurs exactly at the twelve
cap corners. A line-state invariant excludes the six negative regions;
explicit tube zigzags establish the positive residues.

Actual-path existence is eventually periodic for all positive natural
starts and nonnegative natural slopes in $p,r,s$. When $r$ is eventually
at least two, the classification reduces the problem to affine
divisibility. On a ray with $r=1$, a five-chart affine signature has
75 gluing records and twelve marked corners. Its quotient graph is
isomorphic to the tube model. Chart injectivity, degrees, the two-common-
neighbour property and triangle lifting are proved on that model,
supplying the hypotheses of the formal general periodicity theorem.
Residue classes make the twist affine, and their eventual periods combine.
This formal route leaves the separate $r=1$ sufficient conditions above
and their open converse unchanged; it gives no $r=1$ closed classification.

\subsection{Conditional auxiliary lemmas}\label{sec:historical-formal}
\label{lem:historical-kernels}\claimid{lem:historical-kernels}
The auxiliary lemmas have the following hypotheses:
\record{finite_literal_path_decision} assumes a finite type and two
common neighbours; GC lattice transport excludes odd endpoints in a
scaled lattice; the four cylinder arithmetic declarations prove the
height-kernel displacement with stated divisibility hypotheses;
\record{induction_uniform_depth} assumes a finitely branching child
list and genuinely decreasing lexicographic rank. The concrete
\record{ietComponentChildren} models only removal of equal last labels,
not every subtraction, equal-length elimination or accelerated suffix.
The finite decision, scaled-lattice transport and rank statements
are used with their displayed hypotheses.

\label{sec:formal-records}
\section{An extension to fixed oblique directions (proof sketch)}\label{sec:extension}
\label{claim:oblique}\claimid{claim:oblique}
The proposed extension replaces the six side directions by an arbitrary
fixed finite set of primitive lattice directions, keeping eventual
validity, ownership and fixed affine gluing. For a direction $v$, use
transverse coordinate $h=\det(v,z)$. Across a seam the shift is affine.
Let $D$ be a common multiple of the nonzero $|\det(v,e)|$ and put
$Q=4D$. Add the transverse cell modulo $Q$ and $t$ modulo $Q$ to the
finite labels. Entrance and exit indices are
$k_{\rm in}=\lceil\xi_{\rm in}\rceil$,
$k_{\rm out}=\lceil\xi_{\rm out}\rceil$ with
$k_{\rm in}\le k<k_{\rm out}$. Reordering cells $n=\sigma+Qq$
conjugates the discrete midpoint permutation; it does not claim an
order-preserving conjugacy of the continuous intervals.

The return-tower construction keeps a full partition into return towers,
including saved periodic components. Equal final labels remove a
periodic tower; equal lengths merge and remove a label without losing
an orbit. Repeated winner cycles are compressed as powers. The claimed
bound is on the number of Letter/Cat/Power nodes, not on expanded
passage count. Base widths and exponents are affine on fixed residue
branches; word lengths and displacement vectors are rational
polynomials in $t$. A lexicographic measure uses leading coefficient
sum in the original parameter, number of labels, and constant length;
it is not silently reused after coefficient-amplifying reparameterization.

For local ownership, close the side directions under rotations. A
unit triangle meets at most
$\sum_e(\max_\Delta\det(e,z)-\min_\Delta\det(e,z)+1)$
candidate macro-edge levels. Embedded noncrossing edges and finitely
many corner displays bound the length of the local gluing word.
A finite number of charts alone does not bound how often they can
be revisited. Entrance-closed, exit-open ownership excludes repeated
representations of one position from genuine quotient collisions.
Keep the full corner class and the degree-six second-point test;
non-IPR graphs can have fewer than 120 legal starts.

Tower return restores midpoint and direction; two returns restore
left/right phase. Nonempty output, initial suffix and terminal prefix
must be connected before using this as a finite first-repeat bound.
Positions are leaf types and ancestor-power repetition indices.
A free index $j$ contributes $j$ times a displacement depending only
on $t$. Products $jt$ or $jt^2$ are allowed coefficients in one-parameter
Presburger arithmetic; products of two freely quantified repetition
indices are not used. Fix one start address $\theta$ throughout:
\[
\begin{split}
\mathrm{Bad}(t,\theta,e)\ \Longleftrightarrow\
\exists\xi,\eta\,[D_\theta(\xi)\wedge D_\theta(\eta)
\wedge\xi\le e\wedge\eta\le e\\
\wedge\neg I_\theta(\xi,\eta)\wedge G_\theta(\xi,\eta)].
\end{split}
\]
Success would be $\exists\theta,e\,[\mathrm{Legal}(t,\theta,e)
\wedge\neg\mathrm{Bad}(t,\theta,e)]$. Legal includes odd endpoint index,
minimum length, degrees and the preceding output position. Allowing
each collision a different start address would be incorrect.

The remaining obligations are: (i) write the first-crossed-edge,
containment and order inequalities of the local geometric transfer, with an
oblique-family example; (ii) verify the affine tower coefficients
for the $Q$-split IET; (iii) keep effective threshold/period computation
dependent on the effective elimination result \cite{QE}; (iv) compare an oblique-family ray's predicted
tail against full graph decisions. These four conditions delimit the proof sketch.
BGW supplies eventual periodicity only after the geometric formula
has been established.

\subsection{Prospective applications only}
C2 and K4 are specified leapfrogs of C1 and K3; the identified O1$(p,0,s)$
subfamily is a specified leapfrog of O2. The local coordinate map
$(x,y)\mapsto(x-y,x+2y)$ has determinant 3. These are prospects through
the extension, not families in Corollary~\ref{cor:catalog}. Leapfrog
is not assumed to preserve absence in general, and the full
three-parameter O1 family is not classified. These prospects do not
cover arbitrary signatures or every degenerate boundary. For the prospective signatures, twice area is $24a^2+48ab+12b^2$ for C2 and $42r^2+36rs$ for K4. Their $a=0$ and $r=0$ cases collapse faces, while $b=0$ and $s=0$ give zero sides. The zero middle coordinate in the identified O1 subfamily denotes a separately identified leapfrog signature, not permission to discard a zero edge in the old signature.

\section{Fixed-width structure, finite checks and novelty}\label{sec:evidence}
\subsection{Belts and a fixed cylinder}\label{prop:cylinder}
\claimid{prop:cylinder}
Developed triangle midpoints lie on a straight line; vertices alternate
on its two neighbouring integer-height rows. Under $z'=Mz+q$, a height
$h(z)=n\cdot z+e$ transforms as $n'=nM^{-1}$,
$e'=e-n'\cdot q$. If every start is covered by a return preserving
position, direction and primitive height, no cone lies inside a band,
and all widths are at least two, the required pair of endpoint rows
cannot both occur. Width one alone is not sufficient: simplicity and
endpoint-neighbour degrees remain necessary. For integer GC$(k,0)$,
$k\ge2$, the odd-point determinant $\pm1$ contradicts membership in
the $k$-multiple lattice; rotations about parent vertices preserve that
scaled lattice.

For the fixed-cylinder argument let circumference vector $W=gu$,
$g=\gcd(|W_1|,|W_2|)$, $h(z)=\det(u,z)$, and insertion vector $D$
have $H=h(D)>0$. Let $L$ be the least common multiple of the nonzero
$|h(v)|$ over the six unit directions. An entrance-state period is
\begin{equation}\label{eq:cylinder}
 Q=2gL,\qquad {Q\over h(v)}(Hv-h(v)D)\in\Z W.
\end{equation}
Directions with $h(v)=0$ loop in the flat core and repeat without
reaching a cone. If the caps and collars have a fixed finite vertex
set $C$, at most $|C|+1$ new cap visits occur before stopping. Overlap
of the same three-vertex display is not a new visit. On a residue
branch the passages have form $a_i(n)+jv_i$, $0\le j\le\ell_i(n)$,
with affine $a_i,\ell_i$ and fixed $v_i,W$. All collisions are
$a_i(n)+jv_i=a_k(n)+qv_k+zW$, with distinct-position and range
conditions. Ordinary Presburger arithmetic then gives periodicity.
$Q$ is an entrance period, not necessarily the final or least period.

The uniform cap argument does not extend to growing width/caps.
The D12$(3,3t+4)$ example has passage word $(BA)^{t+1}$ and first
legal endpoint index $30t^2+88t+59$. Its $2(t+1)$ passages are unbounded.

\subsection{Finite checks and their domains}\label{sec:finite}
\claimid{sec:finite}
An arithmetic comparison evaluates 1,809 values on nine known A3
rays, and fourteen full graph decisions use an independent constructor.
The fourteen graph inputs are known examples and the arithmetic
values form one sample set. Fixed tail coefficients beyond 64 are
checked on those rays. The E10 descriptor audit checks eight charts,
36 inverse gluings, 108 intervals, 156 connections, 120 descriptors
and twelve labels. A local polygon audit covers 71 charts, 342 edges
and 21 single-zero boundaries. The return-tower audit has 6,564
inputs and 67,312 states. The phase audit covers 48 directions,
14,904 entrances, 4,005 intervals (1,013 empty), 5,496 lifts and
40,380 conjugacy cells.

An independent arithmetic implementation checks 378 coefficient
representations and 13,430 pairs, with maximum period 44 and
threshold 69 in those finite records. Passage-output tests cover
7,056 strips, 2,255 nonempty row outputs, 16,212 subintervals,
99,225 ordered row pairs and 509,160 folded-point comparisons.
These finite ranges check the specified algebra and constructors.

Quotient-model mirror tests compare vertex and edge counts, degree
sequences, seams and path decisions. C1 uses 36 pairs $1\le r,s\le6$;
O2 uses 80 pairs $1\le t\le8$, $1\le d\le10$; D14 uses 81 pairs
$1\le r,s\le9$. Statement-reader samples of 25, 36 and 25 pairs,
respectively, are contained in those ranges. Generated degree and
strip formulas are verified universally by their Lean proofs.

\subsection{Novelty and remaining obligations}\label{sec:novelty}
\claimid{sec:novelty}
\paragraph{Immediate known consequences.}
Once a fixed sentence $\psi_\eta$ describes success on $a^n$,
$t=2n+\eta$, MSO/regular-language equivalence gives a finite automaton;
its unary state orbit is eventually a cycle \cite[Section 3.1]{MSO}.
Finite fixed-modulus residue branches preserve eventual periodicity.
Likewise, a fixed-size parametric Presburger description gives the
conclusion directly by \cite[Theorem 1.15]{BGW}. Neither theorem builds
the geometric description: an unbounded list of path vertices is not
a fixed list of quantified variables, and variable-variable products
are not permitted by that Presburger theorem.

For a fixed-width construction $G_n=C[F^n(B)]$ with fixed finite
boundary, literal-path existence has a fixed CMSO$_2$ description.
Quantify a vertex set $X$, an edge set $E'$ and distinct endpoints
$a,b$: $(X,E')$ is connected, has degree one at $a,b$ and degree two
elsewhere; $|X|\ge4$ is even; $a,b$ have ambient degree five and their
path neighbours ambient degree six; every two-edge path in $E'$ has
adjacent outer vertices in the ambient graph. This permits interior
cones and extra chords. The selected subgraph is precisely a literal
path, in both directions. Finite-state recognizability for a fixed
HR signature \cite[Chapter 1, Theorems 1.12 and 1.21(2)]{Courcelle}
then makes iteration of $F$ a finite-state orbit. Thus this fixed-width
existence statement is a short known consequence, not a new principle.
Mere bounded treewidth of an arbitrary indexed sequence is insufficient:
the uniform repeated construction is essential. The explicit $Q$,
cap-visit bound and actual period-calculation procedure require
separate evaluation.

\paragraph{Nearby results and the missing translations.}
Behmaram--Boutillier's fixed-circumference matching transfer matrix
has $2^m$ states \cite[Section 3, Eq.\ (3.1)]{Behmaram}; increasing
width does not leave the same matrix, and its observable is matchings,
not literal paths. Fischer--Makowsky prove periodicity or semilinearity
of model-size spectra for CMSOL classes with bounded width
\cite[Theorems 1.4, 1.7 and 1.10]{Fischer}. Such a spectrum asserts
existence of some model of a size, rather than deciding a specified
$G_t$. Kotek--Makowsky's labelled (bi-)iterative graph families and MSO
polynomials admit finite matrix transfers and $C^2$-finite recurrences
\cite[Definitions 17--20, Lemma 28, Theorems 30 and 32]{Kotek}.
Their bounded-clique-width framework does not automatically include
growing lattice patches; a recurrence for values is not itself
Boolean periodicity. In each case the geometric family and its
parameter require a uniform interpretation satisfying the cited
hypotheses.

The syntax $\pm x\pm y\le c$ is already the octagon domain
\cite[p.~1]{Mine}. Bozga--Iosif--Kone\v{c}n\'y's difference,
octagonal and finite-monoid affine relations have acceleration and
Presburger transitive-closure results
\cite[Definition 1 and Sections 4.2.1--4.2.2]{Bozga}. These principally
iterate a fixed relation and allow linear drift; here the relation
itself varies with $t$, and all previously output surface vertices
must remain distinct. A fixed-dimensional translation preserving that
condition has not been supplied. The term ``zigzag automata'' there
does not identify the subject with fullerene zigzags.

Jonoska--Krajcevski--McColm study counter languages of walks returning
to a unit cell in periodic digraphs \cite{Counter}; Wanke studies
path/cycle algorithms on finite periodic graphs \cite{Wanke}.
The available abstracts and bibliographic records describe these
walk and algorithm problems. Walks or
decidability alone do not preserve literal simplicity or prove
periodicity of the parameter set.

\paragraph{Unary length is essential.}
The standard WS1S encoding of Presburger integers uses their binary
representations \cite[printed pp.~11--12]{Thomas}. Binary regularity
does not imply eventual periodicity of the represented integers:
the regular language $10^*$ represents powers of two. Our reduction
uses $n$ as unary word length, not the binary digits of $t$.
A separate formula for each $t$ would also be insufficient.

\paragraph{Boundary of the contribution.}
Within the literature comparison of the cited theorem statements
and sections, no prior result giving the full claim from the same
geometric hypotheses has been located. The priority comparison remains
open. The
remaining candidate contribution is the uniform, two-way translation
of geometric passages, seam ownership, legal truncation and all
surface-vertex collisions into one fixed MSO sentence, even with
growing widths and passage counts. Fixed-width theory and the known
logical conclusions do not supply that translation by themselves;
we do not prove that a shorter existing reduction is impossible.
The convex-face reduction is the result proved here. The full
comparison with catalogue and zigzag theory and the extension's
geometric transfer obligations remain separate questions.

The two-interval arithmetic kernel is covered by Li--Chen's uniform
bound for continued fractions of ratios of parameter polynomials
\cite[Theorem 27]{LiChen}. For several intervals, finite residue
branches and a uniform compressed node count, the literature comparison has not located an identical general theorem;
the comparison remains open.

K3's variable-modulus behaviour obstructs a single ordinary
multivariable Presburger description in fixed common coordinates.
Moving to canonical representatives or allowing variable divisibility
does not follow from raywise periodicity. These stronger global-definability questions remain separate. General E10 $r=1$ necessity, untranscribed
families, all degenerations, human proof peer review, and a formal
geometric-to-program pipeline are not completed here. A3 minimum,
density and the IPR census belong to the first article. The D12 classification in Section~\ref{sec:d12-classification} is Lean-proved.


% D12-SECTION-BEGIN
\section{A closed classification for the D12 family}\label{sec:d12-classification}
The D12 member of Corollary~\ref{cor:catalog} also admits a closed
condition. We use the positive integer parameters of the catalogue
\cite{Graver}: put $L=r+2s$, so the carbon order is $12L^2$.
One quotient convention takes two chart-tagged copies of
\[
 H_L=\{(x,y)\in\Z^2:-L\le x\le L,\quad 0\le y\le2L,
                            \quad0\le x+y\le2L\}.
\]
In counterclockwise order its corners are
$(0,0),(L,0),(L,L),(0,2L),(-L,2L),(-L,L)$.
Each side is marked at distance $r+s$ from its initial corner.
The alternating corners and marks have cyclic labels
$(11,0,1,\ldots,10)$ in the first chart and
$(0,11,10,\ldots,1)$ in the second. Equally labelled segments
are glued in reverse order, preserving distance. Quotient vertices
are chart sites modulo these identifications; adjacency is a unit
lattice edge in a common chart.

\paragraph{Classification.}\label{thm:d12-classification}
% \claimid{thm:d12-classification}
For this positive-parameter quotient convention, a literal path exists
if and only if
\[
 \boxed{\quad r\mid(s-1)\quad\hbox{or}\quad r\mid(s+1)
                     \quad\hbox{or}\quad |r-3s|=1.\quad}
\]
Equivalently, the condition is
\[
 r=1\quad\hbox{or}\quad(r=2\ \hbox{and }s\hbox{ odd})\quad\hbox{or}\quad
 \bigl(r\ge3\ \hbox{and }(s\equiv\pm1\pmod r\ \hbox{or }|r-3s|=1)\bigr).
\]
Divisibility by one is automatic; modulo two the two residues both
mean odd $s$, and $|2-3s|=1$ adds only $s=1$.
This classification, and eventual periodicity of path existence on
every positive-natural ray, are Lean-proved for the explicit quotient
model above, with no prose step (Section~\ref{sec:d12-lean}). Its
identification with the catalogue family uses the stated coordinate
conventions. Two further proofs follow. They were developed
separately from each other; each is computer-assisted, and its
all-parameter quotient geometry, the semantics of its computations
and its inductions are prose premises. Complete arguments, explicit
strip formulas, certificate inputs and their checkers accompany the
D12 supplement.

\subsection{A formal proof}\label{sec:d12-lean}
The Lean proof works with the quotient model directly. The two charts
are glued along one boundary circle of length $6L$, and a seam vertex
is recorded by its arc length. For all $r,s\ge1$ the endpoints of an
edge have at most two common neighbours and the twelve labelled points
have degree five; for $r,s\ge2$ they are the only vertices of degree
five. A rotation of
both charts by $60^\circ$ and an exchange of the two charts are
automorphisms of the quotient, so every start can be moved to the
label-$0$ point, from which ten first triangles remain.

Necessity concerns $r,s\ge2$, since the condition always holds when
$r=1$ or $s=1$. For each of the ten starts the forced strip is shown to
repeat a vertex before a legal endpoint whenever the condition fails.
Six starts, and the other four when $s\le r+1$, are settled by a finite
cover of the negative parameter region by integer polyhedra; on each
piece the strip has one fixed combinatorial form. For the other four
starts with $s\ge r+2$ the strip winds around the sphere. Writing
$s=qr+c+t$ with $0\le t<r$ and a fixed offset $c$, the strip consists of
an initial block, a number of laps of four blocks each, and a tail. Each
lap is affine in $r$, $s$ and a lap offset, so induction on the number
of laps carries the strip to the tail. The repetition occurs between
the last lap and the tail unless $s\equiv\pm1\pmod r$.

Sufficiency exhibits explicit literal paths: single blocks for $r=1$
and for $s=1$, a fixed number of blocks on $r=3s\pm1$ and on $s=r\pm1$,
and lap-periodic paths for $s=qr\pm1$, $q\ge2$, whose vertices are
shown distinct for every number of laps. The ray statement follows from
the classification by the arithmetic at the end of this section. The
two-common-neighbour property is proved on the model, so this proof
does not use cyclic edge-connectivity. It follows the strip
decomposition of the two proofs below and is therefore not a third
independent derivation; its correctness rests on the kernel check.

\subsection{A symbolic region cover and cylinder induction}
The first route reduces the twelve degree-five starting sites to ten
representative oriented triangles by two affine lattice symmetries.
The forced-prefix rule of Section~\ref{sec:definitions} then follows
each triangle until a legal endpoint or a repeated quotient vertex.
An exact symbolic executor splits integer affine sign conditions;
the earlier event is constant on each terminal region. Repetition
takes precedence at a tie.

For $r\ge3$, finitely many regions cover $r\ge s-1$ and nine wedges
cover the remaining starts through $s\le10r+1$. Their positive leaves
give the exceptional lines $r=3s\pm1$ and modular boundaries.
The unbounded tail $s\ge10r$ has an overlapping cover by cylinders.
For its four residual starts, one lap increments an auxiliary state
parameter by $2r$. A first exit lies in a window of width $2r$;
the two positive exit planes differ by $r$ and reduce to
$s\equiv\pm1\pmod r$. This is induction on an arbitrary number
of laps, rather than a bound on the number of face passages.
Positive exits are compared with every earlier lap, including
same-chart, seam and parallel collisions. The cases $r=2$ have
a finite complement to the tail, and $r=1$ has explicit boundary paths.

This route requires faithful passage rows and seam maps, exhaustive
integer sign splits, sound rational polyhedral tests, earliest-event
and parity checks, and the full-history cylinder argument. In
particular, its polyhedral emptiness shortcut uses pointedness;
the supplied constraints must satisfy that condition. Finite graph
audits support these semantics on their specified inputs and do not
replace the all-parameter geometric and inductive arguments.

\subsection{Explicit strips and rational certificates}
The second route independently constructs the same quotient, after a
translation of one chart, and derives explicit alternating strips
\[
 (f,A,d,e,n):\qquad A+\lfloor j/2\rfloor d+(j\bmod2)e,
                  \qquad 0\le j<n.
\]
Here $f$ is a chart and $n$ counts vertices. Unit differences
$d,e,d-e$ produce consecutive triangular triples. The negative
parameter regions are covered by short repeated prefixes or by
an arbitrary number of strip pairs followed by a terminal repetition.
The short part uses 29 explicit candidates and 46 region--start covers.
The long part advances a strip offset by $r$ and uses quotient
automorphisms to prove the continuation for every repetition count.
An exhibited repetition with no earlier legal endpoint suffices;
it need not be the first repetition globally.

The infeasibility certificates use nonnegative rational weights
$w_i$ for constraints $f_i\ge0$, with
\[
 \sum_i w_i f_i=-c,\qquad c>0,
\]
and all variable coefficients zero. These identities can be checked
without trusting the search for the weights. Exhaustive integer
coverage splits the negation of $f\ge0$ as $-f-1\ge0$.
The geometric interpretation of the constraints and the arbitrary-strip
induction remain part of the prose argument.

For sufficiency, explicit paths have $6s+4$ vertices when $r=1$,
$4r+6$ when $s=1,r\ge2$, and $56s+12$ on $r=3s+1$, $s\ge1$,
or $56s-10$ on $r=3s-1$, $s\ge2$. For $s=qr\pm1$, $r\ge3$,
$q\ge1$, the alternating constructions have
$2q(3r+5s)+2$ or $2q(3r+5s)$ vertices; the positive-sign
construction extends to $r=2$. Containment, endpoints and seam
ownership are checked, and a global simplicity argument separates
all even/odd rows, including arbitrarily separated strips.
The witnesses allow internal degree-five vertices, as does the definition.

The necessity arguments use cyclic five-edge-connectivity of fullerenes
\cite{Connectivity} to exclude nonfacial dual triangles and justify
unique continuation. The explicit positive unit triangles do not need
that external theorem. The two routes have separately developed
geometry and code; the record gives their derivation-independence and
outcome-exposure qualifications. Human proof review and priority
comparison remain open.

\subsection{Relation to the ray theorem}
The D12 closed condition is an additional classification, not an extra
hypothesis of Theorem~\ref{thm:aligned}. The convex-face decomposition
already supplies D12 ray periodicity through that theorem and its
prose correspondence. The classification also gives the ray result by
elementary arithmetic; the Lean proof takes this route. For
$r(t)=at+b$, $s(t)=ct+d$ on an eventually positive integer ray,
positive $a$ makes both divisibility decisions eventually constant:
$r(t)\mid s(t)-\varepsilon$ implies
$r(t)\mid a(d-\varepsilon)-bc$. If this fixed determinant vanishes,
the ratio is $c/a$ and integrality decides the condition. The
exceptional equality $|r(t)-3s(t)|=1$ is eventually constant as well.
For constant $r=m\ge3$, a valid period divides $m/\gcd(m,c)$;
$m=1,2$ have periods dividing one and two. The Lean ray statement starts at
positive natural coordinates with nonnegative slopes; an eventually
positive integer ray is covered by reindexing its positive tail.
% D12-SECTION-END

\appendix
\section{Local polygons and zero boundaries}\label{app:domains}
\claimid{app:domains}
Twice total polygon area is the carbon order only when the quotient
is valid. The local coefficient audit gives
\begin{center}\small
\begin{tabular}{lll}\toprule
Signature & Side factors & Twice total area\\\midrule
E10 & $p,r,r+s$ & $12r^2+2s^2+12rs+24pr+12ps$\\
D12 & $s,r+s$ & $12r^2+48rs+48s^2$\\
D14 & $r,s$ & $8r^2+24rs+12s^2$\\
D1 & $p,r$ & $36p^2+48pr+12r^2$\\
K3 & $r,s$ & $14r^2+12rs$\\
B2 & $p,r$ & $12p^2+24pr+4r^2$\\
C1 & $a,b$ & $8a^2+16ab+4b^2$\\
O2 & $t,d$ & $14t^2+16td+2d^2$\\\bottomrule
\end{tabular}
\end{center}
The 71-chart/342-edge audit establishes local convexity, positive area
and side lengths in its stated positive domain, not global quotient
simplicity. With other parameters positive, E10 $s=0$ and D12 $r=0$
alone remain locally nondegenerate. E10 $p=0$ or $r=0$, K3 $r=0$,
B2 $p=0$ or $r=0$, C1 $a=0$, and O2 $t=0$ collapse some faces.
D12 $s=0$, either D14/D1 zero, K3 $s=0$, C1 $b=0$ and O2 $d=0$
produce zero sides. Removing an IET zero-length label is not the same
as redefining a graph after cone mergers. The omitted leapfrog-domain
data remain in the extension's prospective records, not the main
corollary. A3's $l=0$ and $k=l$ are the separate A1 and A2 boundaries.

\section{Auxiliary formal statements}\label{app:correspondence}
The auxiliary results retain the following mathematical hypotheses.
\begin{quote}\small
A finite graph with two common neighbours per edge, a cone predicate
and a fixed oriented triangle has a first repeat at most the number
of vertices; a literal path exists exactly when a legal endpoint
occurs before that repeat.\par
A B\'ezout witness for primitive $u$ and
$g\mid Q/\mathrm{cylHeight}(u,v)$ imply that the stated shifted
displacement is a multiple of $gu$, with integer-division and
divisibility hypotheses.\par
A finite-list child relation with a strictly decreasing lexicographic
natural rank has a finite uniform depth from each starting state.\par
The concrete equal-final-label removal component has such a depth.
\end{quote}

\section*{Acknowledgements and AI-use statement}
The author used AI systems to assist with literature
organization, mathematical exploration, proof drafting, implementation,
finite verification, formalization and English exposition. The author
is responsible for the mathematical account and its stated proof
scope. AI systems are not listed as authors.

\section*{License and contact}
Copyright 2026 Tatsuhiko Takimoto. The article text, figures and research
data are available under CC BY 4.0; the accompanying code and Lean
sources are under Apache-2.0, with third-party notices retained.
Contact is through GitHub Issues of the public repository; its URL
has one canonical location in the accompanying
\record{REPOSITORY-URL.txt}. No contact email is listed.

\clearpage
\begin{thebibliography}{99}
\bibitem{A3} Tatsuhiko Takimoto,
\emph{A two-stage classification of linear generalized Stone--Wales paths in
Goldberg--Coxeter images of the dodecahedron}, version 8.1 (2026).
{\urlstyle{rm}\url{https://github.com/tkmt-lab/a3-gsw-classification}}.
\bibitem{BBS} A. Bille, V. Buchstaber, and E. Spodarev,
\emph{Some open mathematical problems on fullerenes}, arXiv:2410.19322,
2024; published 2025, doi:10.1021/acs.jcim.4c01997.
\bibitem{Graver} J. E. Graver,
\emph{Catalog of All Fullerenes with Ten or More Symmetries},
DIMACS Series 69 (2005), 167--188; author version with corrections.
\bibitem{Connectivity} K. Kutnar and D. Maru\v{s}i\v{c},
\emph{On Cyclic Edge-Connectivity of Fullerenes}, arXiv:math/0702511v1,
2007. The cyclic 5-edge-connectivity statement is cited there as known.
\bibitem{MSO} D. Mandrioli, D. Martinenghi, A. Morzenti,
M. Pradella, and M. Rossi,
\emph{Lecture Notes on Monadic First- and Second-Order Logic on Strings},
arXiv:2301.10519v1, 2023, Section 3.1, Theorem 3.2.
\bibitem{BGW} T. Bogart, J. Goodrick, and K. Woods,
\emph{Parametric Presburger Arithmetic: Logic, Combinatorics, and
Quasi-polynomial Behavior}, Discrete Analysis (2017), article 4,
doi:10.19086/da.1254, Theorem 1.15.
\bibitem{QE} A. Mansutti and M. R. Starchak,
\emph{One-Parametric Presburger Arithmetic Has Quantifier Elimination},
LIPIcs 345, MFCS 2025, article 72, doi:10.4230/LIPIcs.MFCS.2025.72.
\bibitem{LiChen} N. Li and S. Chen,
\emph{On the Ring of Integer-valued Quasi-polynomials},
arXiv:0706.4288v2 (2007; internal PDF date 2008), Theorem 27.
\bibitem{Courcelle} B. Courcelle and J. Engelfriet,
\emph{Graph Structure and Monadic Second-Order Logic}, Cambridge
University Press, 2012; author-hosted Chapter 1, Theorems 1.12, 1.21.
\bibitem{Behmaram} A. Behmaram and C. Boutillier,
\emph{Asymptotic enumeration of perfect matchings in m-barrel fullerene graphs},
arXiv:1710.05156, Section 3.
\bibitem{Fischer} E. Fischer and J. A. Makowsky,
\emph{On spectra of sentences of Monadic Second Order Logic with counting},
author draft, 22 May 2003; Journal of Symbolic Logic 69 (2004), 617--640.
\bibitem{Kotek} T. Kotek and J. A. Makowsky,
\emph{Recurrence Relations for Graph Polynomials on Bi-iterative
Families of Graphs}, arXiv:1309.4020v1, 2013.
\bibitem{Mine} A. Min\'e, \emph{The Octagon Abstract Domain},
arXiv:cs/0703084v1, 2007, p.~1.
\bibitem{Bozga} M. Bozga, R. Iosif, and F. Kone\v{c}n\'y,
\emph{Fast Acceleration of Ultimately Periodic Relations},
VERIMAG Technical Report TR-2010-3, 2010.
\bibitem{Counter} N. Jonoska, M. Krajcevski, and G. McColm,
\emph{Counter machines and crystallographic structures},
Natural Computing 15 (2016), 97--113,
doi:10.1007/s11047-015-9527-0.
\bibitem{Wanke} E. Wanke, \emph{Paths and cycles in finite periodic graphs},
MFCS 1993, LNCS 711, 751--760,
doi:10.1007/3-540-57182-5\_66.
\bibitem{Thomas} W. Thomas, \emph{Languages, Automata, and Logic},
Report 9607, May 1996; printed pp.~11--12 (PDF pp.~14--15).
\end{thebibliography}
\end{document}
